Simulation and language-conditioned manipulation benchmarks provide the context for EmbodiedStressBench. Task and policy benchmarks such as RLBench, Meta-World, robosuite, ManiSkill, CALVIN, LIBERO, SIMPLER, RoboCasa, VLABench, and RoboTwin evaluate broad manipulation capabilities, long-horizon instructions, or sim-to-real policy behavior \cite{james2020rlbench,yu2020metaworld,zhu2020robosuite,gu2021maniskill,gu2023maniskill2,mees2022calvin,liu2023libero,li2024simpler,nasiriany2024robocasa,zhang2024vlabench,mu2025robotwin}. Dataset and demonstration resources such as YCB, robomimic, MimicGen, BridgeData V2, DROID, and Open X-Embodiment support object interaction, imitation, and cross-robot learning at scale \cite{calli2017ycb,mandlekar2021robomimic,mandlekar2023mimicgen,walke2023bridgedatav2,khazatsky2024droid,openx2023}. These resources are complementary to the present study: they emphasize policy capability, data coverage, or task success, whereas EmbodiedStressBench isolates the language-to-target generation stage and reports oracle gaps and failure labels.

Language-conditioned manipulation methods also motivate the scope of the benchmark. Transporter Networks, CLIPort, PerAct, VIMA, SayCan, PaLM-E, RT-1, RT-2, RT-X-style data aggregation, and Octo demonstrate increasingly capable links between language, vision, action, and embodied policy learning \cite{zeng2021transporter,shridhar2022cliport,shridhar2022peract,jiang2023vima,ahn2022saycan,driess2023palme,brohan2023rt1,brohan2023rt2,openx2023,ghosh2024octo}. EmbodiedStressBench does not compete with these systems as a policy benchmark. Instead, it asks whether the intermediate target generated from a language-conditioned perception stage is geometrically valid under controlled stressors.

The detector bridge is informed by open-vocabulary perception and segmentation work. CLIP, MDETR, OWL-ViT, Detic, GroundingDINO, YOLO-World, SAM, and SAM2 show that natural-language object localization and segmentation can be modular components in embodied systems \cite{radford2021clip,kamath2021mdetr,minderer2022owlvit,zhou2022detic,liu2023groundingdino,cheng2024yoloworld,kirillov2023sam,ravi2024sam2}. The present detector bridge uses this literature as motivation but keeps the benchmark detector-agnostic: metadata-assisted query-aware selection is a controlled comparator, first detection is a stress-sensitive comparator, and GroundingDINO is one learned plug-in whose query sensitivity can be audited.

Finally, the stress-testing design follows the broader principle that aggregate accuracy can hide systematic failures. CheckList-style behavioral testing and common-corruption robustness benchmarks illustrate how controlled perturbations reveal brittle behavior that average performance misses \cite{ribeiro2020checklist,hendrycks2019robustness}. EmbodiedStressBench applies the same diagnostic stance to robotic target generation by varying detection-list ambiguity, visual occlusion, depth corruption, camera pose, and target-displacement tolerance.
